\section{Compact errors and independent prime boxes}\label{sec:boxes}

The approximation of the original tree by its reference tree is already
adequate at nodes with large gap.  At a node with bounded gap, however, a
large relative error in the small-prime sieve need not be impossible.
We first locate a recent edge on which such an error is visible, and then
place the relevant tuples in boxes of independently varying primes.
The width of the boxes will be chosen last; all constants used to sum
their measures must therefore be independent of that width.

For a tuple $d=p_1\cdots p_m$, always written with
$p_1>\cdots>p_m>w$, write
\[
 x_j=x(p_j),\qquad
 r_j=r_0-\sum_{i=1}^j x_i,\qquad
 v_j=\log_w J_{p_1\cdots p_j}=r_j+a_\star-2.
\]
An endpoint means a tuple, together with its length, in the included
tree of Section~\ref{sec:tree}.  Its prefixes are the preceding tuples
on the same branch.  For a nonempty tuple its strict cutoff is
$b_d=x_m$; the root has the weak cutoff $B$.
For a family $\mathcal A$ of endpoints, put
\begin{equation}\label{eq:boxes-harmonic-measure}
 \mathfrak m(\mathcal A)=\sum_{d\in\mathcal A}\frac1d.
\end{equation}
Different lengths are counted separately.  This is the natural measure
because $\mu_d=YV_0/d$.

Here is the error accounting that the construction will serve.  Let
$\mathcal S$ be any prefix-free collection of included lower nodes at
which the original tree is stopped, and let $\mathcal T$ be its retained
nonstopped nodes.  Expanding the reference tree only as far as these
stops and then using the original tree identity gives
\begin{align}
 S_1(B)\ &\ge YV_0 P_\ee(r_0,B)
   +\sum_{d\in\mathcal S}
       \bigl(S_d(b_d)-\mu_dP_\ee(r_d,b_d)\bigr)
   -\sum_{d\in\mathcal T}|N_d-\mu_d| .
 \label{eq:boxes-error-ledger}
\end{align}
All stopped nodes have even length, so their exact-subtree corrections
have a positive sign.  By Lemma~\ref{lem:small-sieve} and the
exponentially weighted tail estimate in
Proposition~\ref{prop:prime-occupation},
\begin{equation}\label{eq:boxes-large-gap-error}
 \lim_{K\to\infty}\limsup_{z\to\infty}
 \frac{B^2}{YV_0}
 \sum_{\substack{d\text{ included}\\r_d>K}}|N_d-\mu_d|=0.
\end{equation}
Also, for each fixed $K$,
\begin{equation}\label{eq:boxes-small-error}
 \sum_{\substack{d\text{ included},\ r_d\le K\\
                   |N_d/\mu_d-1|\le\eps}}
       |N_d-\mu_d|
 \le C_K\eps\,\frac{YV_0}{B^2}.
\end{equation}
These estimates are over the full included tree, so they remain valid
after any stops have been imposed.  We must control the remaining
compact endpoints, at which the relative error exceeds $\eps$.
At all included endpoints Lemma~\ref{lem:small-sieve} also gives
$|N_d/\mu_d-1|\ll_{a_\star}1$.  Thus discarding a family
$\mathcal A$ of compact endpoints costs
$O_{a_\star}(YV_0\mathfrak m(\mathcal A))$ in
\eqref{eq:boxes-error-ledger}.

\subsection{A recent witnessing edge}

\begin{lemma}[Witnessing edge]\label{lem:edge-witness}
Fix $K$ and $0<\eps<1$.  One can choose a fixed
$K_*>K+10$ such that, with
\[
 H_e=\lceil 2K_*+10\rceil,\qquad
 \delta=\frac{\eps}{4H_e},
\]
the following holds for all sufficiently large $z$, uniformly in the
prescribed residue classes.  If an included endpoint $d$ has
$r_d\le K$ and $|N_d/\mu_d-1|>\eps$, then one of its last $H_e$ edges,
with parent $d'$ and child prime $u$, satisfies
\begin{equation}\label{eq:boxes-witness}
 \left|N_{d'u}-\frac{N_{d'}}u\right|
       >\delta\frac{\mu_{d'}}u,
 \qquad
 x(u),\ \log_w J_{d'}\le 2K_*+3,
 \qquad
 \log_w J_{d'u}\ge a_\star.
\end{equation}
\end{lemma}

\begin{proof}
Choose $K_*$ so that Lemma~\ref{lem:small-sieve} gives
$|N_e/\mu_e-1|<\eps/2$ whenever $\log_w J_e\ge K_*$.  Along the
given branch the numbers $v_j$ decrease strictly, by more than one
at each step.  Let $i$ be the first index for which $v_i<K_*$.
This index exists because $r_d\le K$, and it is not the first index:
$v_1\ge\log_wY-B\to\infty$.

Every nonroot included lower node of cutoff $b$ has both $r\ge2b$
and $r\ge b+2$, by \eqref{eq:admission}.  Its odd child of exponent
$x\le b$ consequently has $r_{\rm child}\ge x$ and
$r_{\rm child}\ge2$.  An included even child has
$r_{\rm child}\ge \mathcal H(x)\ge2x$.  Thus at the crossing index $i$,
regardless of its parity,
\[
 x_i\le r_i=v_i+2-a_\star<K_*+2,
 \qquad v_{i-1}=v_i+x_i<2K_*+2.
\]
All included nodes have $v_j\ge a_\star$.  It follows that there are
fewer than $H_e$ edges from the crossing parent to $d$, and each of
them has parent $v\le2K_*+3$, exponent at most $2K_*+3$, and child
$v\ge a_\star$.

Put $e_e=N_e/\mu_e-1$.  For each edge of prime $u$,
\[
 e_{d'u}-e_{d'}
   =\frac{N_{d'u}-N_{d'}/u}{\mu_{d'}/u}.
\]
The error at the crossing parent has absolute value less than
$\eps/2$, whereas the final error has absolute value greater than
$\eps$.  Telescoping and the triangle inequality show that one of
the fewer than $H_e$ differences has absolute value greater than
$\eps/(2H_e)>\delta$.
\end{proof}

\subsection{Bins, boxes, and regular endpoints}

Fix for the moment $0<\xi<1$.  Partition $(w,z]$ into right-closed
bins of equal logarithmic width, using
\[
 n_{\rm bin}=\left\lceil\frac{\log(z/w)}{\log(1+\xi)}\right\rceil
 \quad\text{bins},\qquad
 h=\frac{\log(z/w)}{n_{\rm bin}\log w}
\]
for their width in the exponent coordinate $x=\log_w p$.
Thus $h\asymp\xi/\log w$ for fixed $\xi$, and the ratio of the upper
and lower endpoints of a bin is at most $1+\xi$.  A bin with lower
endpoint $U$ has
\begin{equation}\label{eq:boxes-bin-count}
 n(U)\asymp\frac{\xi U}{\log U}
\end{equation}
primes, uniformly for $w\le U<z$, once $z$ is sufficiently large
depending on the fixed $\xi$.  This follows from
Lemma~\ref{lem:prime-inputs}; the same conclusion holds for the last
bin because its logarithmic width equals that of the others.
The bin's representative exponent is its right endpoint in exponent
coordinates.

A full box fixes a length and the ordered bin labels of all its
positions.  If a bin appears $m_j$ times, choose any $m_j$ distinct
primes in that bin and put them in decreasing order.  Choices in
different bins are independent.  Give each resulting tuple weight
$1/d$, and write $\mathfrak m(\mathcal B)$ for the sum of these
weights over a full box $\mathcal B$.  Equivalently, the probability
measure obtained by dividing by $\mathfrak m(\mathcal B)$ is a
product, over its used bins, of the reciprocal-product measures on
the chosen subsets.  A position has multiplicity one if its bin is
used just once in the whole box.

A \emph{search bin} is a bin whose lower exponent endpoint is at
least $w^{1/4}$.  We use the fixed constant
$\alpha_{\max}=1/100$.  For a constant $\alpha_{\min}>0$ to be chosen below, a bin wholly
contained in $[\alpha_{\min}B,\alpha_{\max}B]$ in exponent coordinates will be
called an isolated-bin candidate.  When at least one such bin is
used, designate the first such position as the \emph{isolated
position}; its prime will be denoted by $p$.

\begin{proposition}[Regular boxes]\label{prop:regular-boxes}
Fix $K$ and $K_*>K+10$.  For every $\vartheta>0$ there are fixed
constants $C_{\rm len}$ and $\alpha_{\min}>0$, followed by constants
$\xi_0>0$ and $C_{\rm th}<\infty$, with the following properties.
They are chosen before $z$ tends to infinity, and $C_{\rm th}$ is
independent of $0<\xi\le\xi_0$.  Put
\begin{equation}\label{eq:boxes-delta}
 \Delta=2C_{\rm len}\xi.
\end{equation}
For each fixed $0<\xi\le\xi_0$ and sufficiently large $z$, there is
a family $\mathcal G$ of regular included endpoints with $r_d\le K$
such that
\begin{equation}\label{eq:boxes-exception}
 \mathfrak m\bigl(\{d\text{ included}:r_d\le K\}\setminus\mathcal G\bigr)
 \le \frac{\vartheta+C_{\rm th}\Delta+o_\xi(1)}{B^2}.
\end{equation}
The term $o_\xi(1)$ tends to zero after all fixed parameters,
including $\xi$, have been chosen.  The statement is uniform in
the prescribed residue classes.  The regular endpoints and their
covering full boxes have these additional properties.

\begin{enumerate}[label=\textup{(\roman*)}]
\item Every regular endpoint has length at most
$C_{\rm len}\log B$ and has an isolated position.  Every used
search bin, including its isolated bin, has multiplicity one.
Every position satisfying
\[
 x_j\le2K_*+3,\qquad r_{j-1}\le2K_*+5
\]
also has multiplicity one.  The isolated position precedes every
edge supplied by Lemma~\ref{lem:edge-witness}.

\item For every nonempty even prefix of a regular endpoint,
\begin{equation}\label{eq:boxes-threshold-clearance}
 r_j-\mathcal H(x_j)>10\Delta.
\end{equation}
Moving prime choices anywhere within their fixed bins changes any
prefix gap by at most $\Delta$, and changes a threshold expression
$r_j-\mathcal H(x_j)$ by at most $\Delta+2h$.  In particular the
representative gaps and exponents of every regular box satisfy
\begin{equation}\label{eq:boxes-representative-safety}
 r_j^{\rm rep}\ge \mathcal H(x_j^{\rm rep})+3\Delta
 \quad\text{at every nonempty even prefix}.
\end{equation}

\item Let $\mathfrak B$ be the distinct full boxes containing at
least one regular endpoint.  There is a constant $C_K^{\rm box}$,
depending only on a fixed enlargement of the endpoint compact
range and the allowed starting ratios, such that
\begin{equation}\label{eq:boxes-expanded-mass}
 \sum_{\mathcal B\in\mathfrak B}\mathfrak m(\mathcal B)
       \le\frac{C_K^{\rm box}}{B^2}.
\end{equation}
In particular this constant is independent of $\xi$, $\vartheta$, and
the total number of boxes.  Every tuple in these boxes is a
standard path from initial gap $r_0+1$, ends at gap at most $K+2$
for that starting gap, and has its actual final
$\log_wJ\ge3a_\star/4$.

\item In a regular box, declare an edge a candidate if it is
among the last $H_e=\lceil2K_*+10\rceil$ positions, is after the
isolated position, has multiplicity one, and can satisfy the size
bounds in \eqref{eq:boxes-witness} for some tuple in the box.
There are at most $H_e$ candidates, and every witness for a regular
endpoint is among them.  Throughout the whole box of each candidate,
writing its parent as $d'=pq$ and its prime as $u$, one has
\begin{equation}\label{eq:boxes-candidate-slack}
 \log_w\frac{Y}{pq}\le2K_*+4,
 \qquad
 \log_w\frac{Y}{pqu}\ge\frac{3a_\star}{4}.
\end{equation}
The bins for $p$ and $u$ have lower endpoints $R$ and $U$ satisfying
\begin{equation}\label{eq:boxes-prime-ranges}
 z^{\alpha_{\min}}\le R\le z^{\alpha_{\max}},
 \qquad w\le U\le w^{2K_*+4}.
\end{equation}
The choices of $p$, $q$, $u$, and the descendants after this edge
are independent under the full-box measure.  All their prime
factors are disjoint.  If $S_q$ is the product of the bin upper
endpoints for the factors of $q$, with their specified
multiplicities, then
\begin{equation}\label{eq:boxes-sq}
 1\le\frac{S_q}{q}\le(1+\xi)^{C_{\rm len}\log B}.
\end{equation}

\item Every regular box has at least $c_7\log w$ search
positions, for an absolute $c_7>0$ and sufficiently large $z$.
These positions have distinct bins.  All its nonsearch positions
follow all its search positions.
\end{enumerate}
\end{proposition}

\begin{proof}
We first establish the exclusions leading to \eqref{eq:boxes-exception}.
The compact-history estimate in
Proposition~\ref{prop:prime-occupation} permits each geometric
exception to be estimated in the tilted prime chain, at a cost at
most $C_KB^{-2}$ times its occurrence probability.  Histories with
a ratio exceeding $S_*=(\log B)^2$ cost $O_A(B^{-A})$ in the
original measure, by Lemma~\ref{lem:high-states}.  On the remaining
histories, Proposition~\ref{prop:prime-coupling} compares the prime
chain with the continuous chain up to
\[
 N_*=O((\log B)^3)
\]
steps.  Its failure probability tends to zero, and its matched
states satisfy $|s-\widetilde s|\le h_0$ and
$|\log(r/\widetilde r)|\ll N_*h_0$, with $h_0$ smaller than every
fixed negative power of $\log B$.  These errors remain negligible
in all polynomially many opportunities below.  We work with the
continuous comparison where useful, retaining the indicated slack
when returning to prime paths.

\paragraph{Length.}
The iid regeneration cycles from Lemma~\ref{lem:regeneration} have
finite mean length and positive finite mean log decrement. Their
exponential moments imply finite variances, so Chebyshev's inequality
applied to the iid cycle sums gives their respective weak laws.
Inverting the positive mean decrement and bounding the last incomplete
cycle then gives a positive limiting log decrement per step in
probability. Indeed, the largest length among $O(\log B)$ cycles is
$o(\log B)$ in probability by the exponential moment. The initial
segment is uniformly tight for the allowed starting ratios.
Choose a fixed $C_{\rm len}$ larger than twice the reciprocal of
this speed.  The probability that more than $C_{\rm len}\log B$
steps occur before the accumulated decrement exceeds
$\log r_0+O(1)$ tends to zero.  An included compact endpoint has
$r_d\ge2$, so its accumulated decrement is at most that quantity.
The coupling transfers this conclusion.  The excluded harmonic
mass is $o(B^{-2})$.

\paragraph{An isolated position.}
Put $\theta=\alpha_{\max}/4$.  A path from the initial gap to a
fixed compact must cross below $\theta B$.  After the first step,
each transition of a standard path leaves at least half of the
old gap: an even parent has next exponent at most half its gap,
and an odd parent with an admitted even child has next exponent
at most one third of its gap.  The first step cannot itself make
this crossing when $z$ is large.  Thus the crossing arrival has
\[
 .4\theta B\le r\le\theta B.
\]
Consider the slightly enlarged gap band
$[.3\theta B,1.1\theta B]$.  By Lemma~\ref{lem:band-bounds}, the
probability of any arrived-at state in this band with ratio
$s>T$ tends to zero uniformly in $B$ as $T\to\infty$.
Indeed the expected count is bounded by a constant times the
tail integrals of the invariant densities, together with the
corresponding first-step density; all these tails tend to zero.

Choose $T$ so that this probability costs less than $\vartheta/2$ in
units of $B^{-2}$.  At the crossing arrival with $s\le T$, its last
exponent is $x=r/s$, so
\[
 \frac{.4\theta}{T}B\le x\le\theta B.
\]
Choose $0<\alpha_{\min}<.2\theta/T$.  The entire bin of this
factor then lies in $[\alpha_{\min}B,\alpha_{\max}B]$ for large
$z$, also after the coupling errors.  Designate the first such
used bin.  Since $\alpha_{\min}B\to\infty$, its position precedes
every possible witnessing edge of
Lemma~\ref{lem:edge-witness}.  Also
$\alpha_{\min}B>w^{1/4}+h$ eventually, so its bin is a search bin.

\paragraph{Repeated search bins.}
If two positions have the same bin label, some two adjacent
positions have that label.  Write $x'$ for the earlier exponent,
$x$ for the next one, $s=r/x'$ for the ratio at their parent, and
$t=r/x-1$ for the next ratio.  Then
\begin{equation}\label{eq:boxes-repeat-width}
 0\le t-(s-1)=r\left(\frac1x-\frac1{x'}\right)
       \le\frac{s h}{x}.
\end{equation}
For a search bin, $x\ge w^{1/4}$.  On low-state histories, the
continuous comparison must therefore draw into a neighborhood of
its lower support endpoint of width at most
\[
 O\left((S_*+1)
       \left(\frac{h}{w^{1/4}}+N_*h_0\right)\right).
\]
The densities of the kernels in \eqref{eq:tilted-kernel} on such
pieces are bounded by a polynomial in $S_*$, as proved in the
prime-coupling argument.  Summing this conditional probability
over at most $N_*$ draws gives $o(1)$.  Here every polynomial in
$\log B$ is negligible compared with $w^{1/4}$, and the same
polynomials times $h_0$ tend to zero.  We discard all these repeats.

\paragraph{Repeated bins at a possible terminal edge.}
Next discard any repetition involving a position with
$x_j\le2K_*+3$ and $r_{j-1}\le2K_*+5$.  It suffices to compare
with an adjacent position in the same bin.  If that neighbor is
the preceding position, the comparing draw is the edge at $j$;
if it is the following position, the comparing draw is at $j+1$.
In either case the parent gap is already in a fixed compact.
Since all used exponents exceed 1, its ratios and those of the
comparing draw are bounded in terms of $K_*$.  Formula
\eqref{eq:boxes-repeat-width} now gives a width $O_{K_*}(h)$,
enlarged under coupling to $O_{K_*}(h+N_*h_0)$.
Only boundedly many draws can have parent gap in this compact
before cutoff exit, because each such step decreases the gap by
more than 1.  The kernel densities are bounded there.  The
probability of this exception is consequently $o_\xi(1)$.
This verifies all multiplicity assertions in part~(i), and every
witness has the required singleton bin because its parent
$r\le2K_*+5$.

\paragraph{Near-threshold even prefixes.}
Discard paths having a nonempty even prefix for which
$0\le r_j-\mathcal H(x_j)\le10\Delta$.  This event is estimated before
imposing the length restriction, so it suffices to bound its
probability by $O(\Delta)+o(1)$.
At its preceding odd state, let the gap be $R$ and the next ratio
be $t$.  Then $x=R/(t+1)$ and $r_j=Rt/(t+1)$.  If $x\ge2$,
\[
 r_j-\mathcal H(x)=\frac{R(t-2)}{t+1},
\]
so the allowed ratios lie within $O(\Delta/R)$ of 2.  If $1<x\le2$,
\[
 r_j-\mathcal H(x)=\frac{R(t-1)}{t+1}-2.
\]
For small $\Delta$ this confines $t$ to a fixed neighborhood of
$[2,3]$, and differentiation in $t$ again gives a derivative
comparable to $R$.  The allowed width is $O(\Delta/R)$ in this
case as well.  Both descriptions are valid with one-sided
endpoints when the two branches meet at $x=2$.

At any such draw the preceding ratio is at most $t+1$, and hence
is bounded by an absolute constant.  The kernel density is
therefore bounded.  In a unit band of $\log R$, the number of
preceding states having this bounded ratio is deterministically
bounded: each such arrived-at state has a log decrement
$g(s)=\log(1+1/s)$ bounded below by a positive constant.  Thus
the conditional widths sum over the bands to
\[
 O\left(\Delta\sum_{j\ge0}e^{-j}\right)=O(\Delta),
\]
since all possible parent gaps are bounded away from zero.
Coupling enlarges each ratio width by $O(N_*h_0)$; summing over
the $O(\log B)$ relevant bands gives $o(1)$.  This proves the
claimed near-threshold probability bound.

Combining the exclusions with the compact-history bound proves
\eqref{eq:boxes-exception}, after increasing $C_{\rm th}$ and
using the spare part of $\vartheta$ to absorb any fixed adjustable
probability costs.  All density, band, and compact-history
constants were fixed independently of $\xi$.

\paragraph{Movement and representative safety.}
For a tuple of the retained length,
\[
 m h\le C_{\rm len}\log B\,
           \frac{\log(1+\xi)}{\log w}
       \le2C_{\rm len}\xi=\Delta
\]
for sufficiently large $z$, since $\log B/\log w\to1$.
This bounds the movement of every prefix gap.  The function $\mathcal H$
is 2-Lipschitz, proving the threshold movement bound in part~(ii).
Choose $\xi_0$ small enough that
\begin{equation}\label{eq:boxes-width-requirements}
 \Delta<\min\{a_\star/4,1/4\}
 \quad\text{and}\quad \Delta+2h<1
\end{equation}
for all sufficiently large $z$.  Since also $h=o_\xi(\Delta)$,
the clearance in \eqref{eq:boxes-threshold-clearance} leaves more
than $3\Delta$ at every representative even prefix.  This proves
\eqref{eq:boxes-representative-safety}.  In particular these
representative checks depend only on the labels of the box.

\paragraph{The measure of the full-box enlargement.}
Take a box containing an original regular endpoint, and let an
arbitrary tuple in that full box be given.  At every even prefix
the original tuple satisfies $r_j\ge \mathcal H(x_j)\ge2x_j$.  Movement
to the new tuple can lose at most $\Delta+2h$ in the expression
$r_j-2x_j$.  Increasing the initial gap from $r_0$ to $r_0+1$
therefore makes every even prefix a valid standard admission.
The odd steps impose no extra standard admission test.  The
endpoint gap for this shifted start is at most
$K+\Delta+1<K+2$.

The actual, unshifted endpoint has $v\ge a_\star-\Delta
\ge3a_\star/4$.  Thus all enlargements also retain the positive
length margin needed later for interval freezing.

Distinct boxes of a given length have disjoint tuple sets.
Although different lengths are counted in the sum of box measures,
Proposition~\ref{prop:prime-occupation} bounds precisely the sum
of original weights over all standard endpoints in a fixed compact.
Apply it from $r_0+1$ with ending gap at most $K+2$.  This start has
ratio in $[1.99,2.3]$ for large $z$, just as required by that
proposition.  We obtain \eqref{eq:boxes-expanded-mass} with a
constant depending only on this fixed compact and the starting
range.  No counting factor for the boxes occurs, and this
constant is independent of their width or their defining
regularity parameters.

\paragraph{Candidate edges and independence.}
Every witness of a regular endpoint is among its last $H_e$
positions, lies after its isolated position, and has a singleton
bin, by the preceding arguments.  Hence it is a candidate.
If a candidate can satisfy the witness size bounds at one tuple,
the gap movement at another tuple is at most $\Delta$.  The first
requirement in \eqref{eq:boxes-width-requirements} and $\Delta<1$
give \eqref{eq:boxes-candidate-slack}.  The corresponding exponent
movement is at most $h<1$, giving the range for $U$ in
\eqref{eq:boxes-prime-ranges}; the range for $R$ follows from the
definition of the isolated position.

Both the isolated bin and the candidate bin have multiplicity
one.  Moreover, because the labels are ordered and the candidate
bin occurs only once, no bin can contain both a factor before
that edge and a factor after it.  Removing the isolated factor
from the parent therefore partitions the remaining choices into
disjoint sets of bins for $q$, for $u$, and for its descendants.
Their normalized measures are independent by the product
description of a full box.  They are also independent of the
isolated prime.  Sorting and distinctness within any other bin
do not change this conclusion.  Inequality~\eqref{eq:boxes-sq}
follows by multiplying the ratio bound $1+\xi$ at every factor
of $q$.

\paragraph{The number of search positions.}
The total exponent contributed by nonsearch positions is at most
\begin{equation}\label{eq:boxes-nonsearch-sum}
 C_{\rm len}\log B\,(w^{1/4}+h)=o(w^{1/2}).
\end{equation}
They all occur after the search positions, by decreasing bin
order.  In one regular tuple in the box, let $m_s$ be the number
of search positions.  Immediately after its last search factor
the gap is at most $K+o(w^{1/2})\le w^{1/2}$ for large $z$.
Each standard transition retains at least one third of its
preceding gap, including the first step for sufficiently large
$z$.  Since $r_0\ge B$ eventually, this implies
\[
 3^{-m_s}B\le w^{1/2},\qquad
 m_s\ge\frac{\log B-\tfrac12\log w}{\log3}.
\]
Now $\log B/\log w\to1$, so the last expression is at least
$c_7\log w$ for an absolute positive $c_7$, for example
$c_7=1/10$, once $z$ is sufficiently large.  The search labels
are fixed throughout the full box, and their multiplicities
are one.  This proves part~(v) and completes the proposition.
\end{proof}

\begin{remark}\label{rem:boxes-order-and-sums}
Proposition~\ref{prop:regular-boxes} separates the adjustable
exceptions from the final width choice.  After $K$ and $K_*$ are
fixed, one first chooses $\vartheta$, $C_{\rm len}$, and
$\alpha_{\min}$.  The bounds for full-box mass and for compact
history events are then fixed independently of $\xi$.
Consequently a later statement failing on a fraction at most
$\sigma$ of each full candidate box costs at most
\[
 \sigma H_e\sum_{\mathcal B\in\mathfrak B}
                      \mathfrak m(\mathcal B)
 \le\sigma H_e C_K^{\rm box}B^{-2}.
\]
This remains true if the number of boxes grows with $z$ or $\xi$.
For a singleton bin, count fractions and harmonic fractions
differ by a factor at most $1+\xi\le2$, so bounds conditional on
all other factors and measured against all primes in that bin
have the same summed interpretation.  Finally $\xi$ may be
chosen sufficiently small to make $C_{\rm th}\Delta$ and all
required movement errors small.  The prime-count asymptotics
are only used after this fixed choice, so their eventual
thresholds may depend on $\xi$ without changing the constant
order.
\end{remark}
